% ============================================================
% Seccion: Plan de trabajo hasta el final del semestre
% Los porcentajes de avance son una estimacion del grupo: ajustarlos aqui
% ============================================================
\section{Plan de trabajo}

% ============================================================
\begin{frame}
\frametitle{Estado de avance y plan de trabajo}
\vspace{-0.5em}
\begin{columns}[T]
\begin{column}{0.42\textwidth}
{\small \textbf{Avance por etapa}}
\vspace{0.5em}

\begin{tikzpicture}[x=1cm, y=1cm]
  \foreach \y/\nombre/\p in {%
      0/{1. Exploración de los datos}/100,
      -1.0/{2. Visualización con autoencoder}/90,
      -2.0/{3. Clustering: GMM y MCDSVDD}/80,
      -3.0/{4. Curvas de luz y outliers}/60}{
    \node[font=\footnotesize, anchor=west, inner sep=0pt] at (0,\y+0.32) {\nombre};
    \fill[structure!18, rounded corners=1pt] (0,\y-0.12) rectangle (4.4,\y+0.08);
    \fill[structure, rounded corners=1pt] (0,\y-0.12) rectangle ({4.4*\p/100},\y+0.08);
    \node[font=\footnotesize\bfseries, text=structure, anchor=west] at (4.5,\y-0.02) {\p\,\%};
  }
\end{tikzpicture}

\vspace{0.6em}
{\footnotesize \textbf{Resultado esperado}: un ranking de anomalías depurado de errores de datos,
evaluado con AUROC y respaldado con las curvas de luz.}
\end{column}

\begin{column}{0.55\textwidth}
{\small \textbf{Próximos pasos}}
\vspace{0.1em}
{\footnotesize
\begin{enumerate}\itemsep=3pt
\item Experimento de anomalías: ocultar una clase, reentrenar y medir si el modelo la marca
  como rara (AUROC)
\item Filtrar períodos en alias y objetos con pocas detecciones antes de calcular el puntaje
  de anomalía
\item Completar la etapa 4 en MCDSVDD: pertenencia a más de un cluster y sub-clusters (ya
  hecho para el GMM)
\item Ajustar los hiperparámetros de MCDSVDD
\end{enumerate}
}
\vspace{0.2em}
{\small \textbf{Aristas posibles}}
\vspace{0.1em}
{\footnotesize
\begin{itemize}\itemsep=2pt
\item Autoencoder solo con período y color para el GMM: en validación el ARI sube de 0,21 a
  0,33
\item Latente distinto por modelo: con 32 dimensiones el GMM también llega a 0,33
\item Entrenamiento conjunto de autoencoder y GMM, como en DAGMM
\end{itemize}
}
\end{column}
\end{columns}
\end{frame}
